\documentclass[10pt,a4paper]{amsart}
\usepackage[margin=19mm]{geometry}
\usepackage{amsmath,amssymb,mathtools}
\usepackage[scr=rsfs,cal=euler]{mathalfa}
\usepackage{xfrac}
\usepackage[unicode,hidelinks,bookmarks=false]{hyperref}
\newcommand{\ie}{i.e.\ }
\newcommand{\cf}{cf.\ }
\newcommand{\resp}{respectively\ }
\newcommand{\Subsection}{\subsection}
\providecommand{\nobreakdash}{\nobreak\hbox{-}}
\setlength{\emergencystretch}{2em}
\multlinegap0pt
\usepackage{bm}
\usepackage{bbm}
\usepackage{dsfont}
\usepackage{xspace}
\usepackage{cancel}
\usepackage{mathtools}
\usepackage{amsthm}
\makeatletter
\@ifundefined{theorem}{\newtheorem{theorem}{Theorem}}{}
\@ifundefined{corollary}{\newtheorem{corollary}[theorem]{Corollary}}{}
\@ifundefined{lemma}{\newtheorem{lemma}[theorem]{Lemma}}{}
\makeatother
\DeclareMathOperator{\Bell}{Bell}
\title[Typical $T$-free graphs]{Typical $T$-free graphs}
\author{Bruce Reed, Yelena Yuditsky}
\date{Source: arXiv:2506.01067v1 (CC BY 4.0)}
\begin{document}
\maketitle

\noindent\emph{Source and licence.} Typical $T$-free graphs. Version arXiv:2506.01067v1 (CC BY 4.0), distributed under the Creative Commons Attribution 4.0 International licence, as recorded in the arXiv licence metadata for this exact version and shown on the abstract page https://arxiv.org/abs/2506.01067v1. Full rights notice: licenses/arxiv-2506.01067-CC-BY-4.0.txt.\par\smallskip

\noindent\textit{Editorial scope.} Counted unit: the Lemma whose source label is \texttt{nobigstable} in the article (source0.tex lines 1595--1602, source label \texttt{nobigstable}), delivered together with the article's own complete original proof of it (source0.tex lines 1604--1766, one \texttt{proof} environment of 163 lines). No mathematics is written, repaired or re-derived: statement and proof are the article's own text line for line. Required context retained uncounted: dangerous2cor (lines 1203--1207); good (lines 1292--1300). Every definition, display and label of those lines is retained unchanged. Omitted from this package and not redistributed: 1--1202, 1208--1291, 1301--1594, 1603--1603, 1767--1881. Nothing inside a retained excerpt is omitted or altered: every retained body is delivered exactly as the article's lines are written, so no omission receipt of any kind accompanies this package. Citations: the retained excerpts carry no citation command at all, so no bibliography entry of the article is transcribed and no reference is redistributed. Reference adaptation: none was needed and none was made; every reference inside the delivered unit resolves inside it, so no unresolved reference or citation is produced. Printed slips kept literal: no printed word, symbol, hypothesis or number of the counted statement or of its proof was altered. Layout and class: the article's own class is \texttt{article} with packages including amsthm, graphicx, caption, subcaption, amssymb, amsmath, xfrac, color, hyperref, commath, mathrsfs, ulem, cleveref, todonotes; the wrapper is the lane's pinned \texttt{amsart} editorial wrapper at 10pt on A4, which restates the theorem-like environments the retained text needs and adds the article's own macro definitions that the retained text needs (\texttt{\textbackslash Bell}), with extra wrapper packages (bm, bbm, dsfont, xspace, cancel, mathtools). The delivered numbering is the unit's own. Assets: no retained span contains a figure environment, a graphics-inclusion command, a PDF-page inclusion, a table or a TikZ picture; no image file is copied into this package, the package declares no asset dependency and no image file is redistributed with it. Licence: version arXiv:2506.01067v1 of this article is distributed under the Creative Commons Attribution 4.0 International licence, as recorded in the arXiv licence metadata for this exact version and shown on the abstract page https://arxiv.org/abs/2506.01067v1. A full rights notice is delivered at licenses/arxiv-2506.01067-CC-BY-4.0.txt. Characters: every retained excerpt is pure ASCII, so the delivered engine, XeLaTeX with its default Latin Modern text font, reports no missing character.
\begin{lemma}\label{dangerous2cor}
Suppose, for some  graph $H$ and some $q$ and $w$,  we are given a partition $\pi=(\pi_1,\pi_2,\ldots,\pi_w)$ of $[n]$ and a pattern $(F_1,F_2,\ldots,F_w)$ on $\pi$. Moreover we are given a partition of $H$ into $S_1,S_2,\ldots, S_w$ and a set $D\subset \pi_k$ for some $k\in [w]$. 
Then, the  total number of extensions of this pattern to an $H$-free graph  such that the choices for the edges from $D$ 
to $V-\pi_i$ are $q$-choice destroying   is at most $2^{m_\pi-\gamma  q^2}$ for some $\gamma>0$ which depends only on $H$. 
\end{lemma}

\begin{corollary}\label{good}
For sufficiently small $\epsilon>0$ there is a constant $C'_T\in \mathbb{N}$ such that almost every $T$-free graph $G$ on $[n]$ extends a really $\epsilon$-relevant  pattern on some partition $\pi$ of $[n]$ such that the following holds: 
\begin{itemize}
\item [(i)]  for  any coarsening of  $\pi$ and set $D$ in some part of the coarsening,  the choices of the edges out of $D$ are not  $n^\frac{2}{3}$-choice destroying set,  and
\item [(ii)] there is a set $H$ of size  at most $C'_T \log n$ consisting of the union of   a maximal family of   $\frac{n}{100(\alpha(T)-1)^2}$-dangerous set 
of size at most 8 in the coarsening, and hence intersecting  every $\frac{n}{100(\alpha(T)-1)^2}$-dangerous set  of size at most 8 in the coarsening. Furthermore for some 
$\mu>0$ which depends only on $T$ and $\epsilon$ there are $2^{O(n)+n|H|(1-\frac{1}{\alpha(T)-1}-\mu)}$ choices for $H$ and the edges from it. 
\end{itemize}
\end{corollary}

\begin{lemma}\label{nobigstable}
The  bad $T$-free  graphs extending an $\epsilon$-relevant pattern $F_1,...,F_{\alpha(T)-1}$  over some partition $\pi$ such that we can choose $H$ so that  (i) and (ii) of Lemma \ref{good} hold,  there is no   $i^*$    such that $\pi_{i^*}$  contains  a stable set of size exceeding $ \frac{6n}{10(\alpha(T)-1)}$, there is no $F_i-H$ which is very far from a clique  
 and  there is an  $i^*$ satisfying the following are   an $o(1)$ proportion of the $T$-free graphs on $n$ vertices:
 
There is a subgraph $C$  of $F_{i^*}-H$  such that   $C$ has size at least $\frac{n}{20(\alpha(T)-1)^2}$, 
every  component of $\overline{C}$ has size  at most  $n^{1-\frac{1}{5(\alpha(T)-1)}}$  and is the disjoint union of a vertex and a clique of size at least two.
Furthermore, the components of $\overline{C}$ lie in different components of $\overline{F_{i^*}-H}$.
 \end{lemma}

\begin{proof}

We define a set $H'$ iteratively as follows. Initially $H'$ is empty. 
For each component $J$ of $C-H' $ we let $v_J$ be a vertex which in $G$ sees none of $J-v_J$. 
For $v\in V-F_{i^*}-H-H'$ we define $J'(v)$ to be the set of vertices of $J$ which are adjacent to $v$ if $v_J$ is not, or are nonadjacent to $v$ otherwise. 

To obtain  $H'$ we  repeatedly add to it  if possible either 
\begin{enumerate}
\item[(a)]  a vertex which  is  in $F_i-H-H'$ for some $i \neq {i^*}$  and such that either $v$ is nonadajacent to fewer than $\frac{|C-H'|}{10}$
vertices of $C-H'$ or the sum over all the components  of  $\overline{C-H'} $ of  $|J'(v)|$
is less than $\frac{|C-H'|}{10}$, or 
\item[(b)] a triple of vertices   in $F_{i^*}-H-H'$  such that  for some $i \neq {i^*}$  there is a vertex $u$ of $F_i-H-H'$ such that $u$  has at least
$\frac{n}{10(\alpha(T)-1)}$ neighbours in $F_i-H-H'$  and all of these of have an edge to the triple.  
\end{enumerate}

We note that for some $\psi>0$ depending only on $T$,  the number of choices for the vertices of $H'$ and the edges from them is $O(2^{bn}2^{|H'|(1-\frac{1}{\alpha(T)-1}-\psi)n})$. 
So, for all $T$ there is a $C'_T$ such that  the subset of the $T$-free graphs  we are considering for which $|H'|>C'_T \log n$ is a $o(1)$ proportion of the $T$-free graphs. 
Hence we can  and do restrict our attention to $G$ for which we stop growing $H'$ before it contains $C'_T \log n$ vertices. 


We note that at this point, for every $v \in V-F_{i^*}-H-H'$  there  is a set $T_v$ of at least $\frac{n^{\frac{1}{6\alpha(T)-1}}}{2000(\alpha(T)-1)^2}$ disjoint triples  in $C$ which together with $v$ induce a $P_4$, in which $v$ is an endpoint. 

Suppose $T$ is spiked. 

The maximality of $H$,  implies  none of these $P_4$s extend to an  $M$ in $F_{i^*} \cup F_i-H$.
For each $i \neq i^*$,  for each component $K$ of $F_i-H-H'$ which has size exceeding 2
we apply  this fact to a vertex  $v$ in $K$ missing the rest of $K$. 


We let $P_v$  be a family of $\lfloor \frac{|K|-1}{2} \rfloor>\frac{|K|}{3}$ pairs 
of vertices in $K-v$. For each pair in $P_v$ and triple in $T_v$ there is a choice of edges between them which yields that together 
with $v$ they yield an $M$. Thus, the number of choices of edges between the union of the elements of $T_v$ and the union of the elements 
of the $P_v$ is at most $2^{6|T_v||P_v|}(1-\frac{1}{2^6})^{|T_v||P_v|}= O( 2^{6|T_v||P_v|- n^{\frac{1}{7(\alpha(T)-1)}}|K|})$. 
On the other hand,  since $K-v$ is a complete multipartite,  the number of choices for $K$ and the edges from its vertices within $F_i-H-H'$  is at most ${n \choose k}n \Bell(|K|-1)
=O(2^{(2|K|+1)\log n})$. 

We consider next the  number $b$  of pairs of vertices forming  components of $F_i-H-H'$. We know  we have $b$ vertex disjoint $P_3$ each contained in $F_i$ for some $i \neq i^*$. Now, there are $n^{\frac{1}{6(\alpha(T)-1)}}$
disjoint $\overline {P_3}$ in $C$. Hence since $H$ is maximal the number of choices for edges between the elements of  $\pi$ 
is at most $ 2^{m_\pi}(1-\frac{1}{2^9})^{bn^{\frac{1}{7(\alpha(T)-1)}}}$.
It follows that the $T$-free graphs  extending a pattern we are considering extends one for which $b \ge n^{1-\frac{1}{3(\alpha(T)-1)}}$ 
is a $o(1)$-proportion of the $T$-free graphs, 
and we can focus on $G$ for which this is not  the case.
But then, since we can choose the two vertex components of the $F_i$ for $i \neq i^*$  in at most $n^{2b}$ ways, it follows that given  the choice for the pattern 
on $F_{i^*}-H-H'$, there are $2^{O(n)}$ choices for the rest of the pattern on $G-H-H'$. But by the maximality of $H$,  the components  of $\overline{F_{i^*}-H-H'}$
are the disjoint union of a vertex and a complete multipartite graph. Hence the number of $T$-free graphs extending the patterns we are considering 
is a $o(1)$ proportion of the $T$-free graphs certified by partitions into two complements of a matching and $\alpha(T)-3$ cliques.

So $T$ is not spiked. 

By the maximality of $H$,  for every $v \in F_i-H-H'$, none of the $P_4$s  induced by $v$ and a triple in $T_v$ extend to an  $P_6$ in $F_{i^*} \cup F_i-H$.
For each $i \neq i^*$,  for each component $K$ of $F_i-H-H'$ which has size at least 2. 
We apply  this fact to a vertex  $v$ in $K$ missing the rest of $K$. 


If $|K| >20$, we greedily  select a family $P_v$   of $\lceil \frac{|K|}{20} \rceil$ pairs 
preferring pairs of adjacent vertices. For each  adjacent pair in $P_v$ and triple in $T_v$ there is a choice of edges between them which yields that together 
with $v$ they yield a $P_6$. If $P$ is a pair of nonadjacent vertices in $P_v$ then there is a stable set in $K$ containing all but  less than $\frac{|K|+1}{10}$ of its 
vertices. Hence $|K|<\frac{4n}{5(\alpha(T)-1)}$. Thus there are at least $\frac{n}{10(\alpha(T)-1)}$ vertices in $F_i-H-H'-K$ and this set which forms the neighbourhood of $v$ 
in $F_i-H-H'$ is also the common neighbourhood of $v$ and the vertices in $P$ in $F_i-H-H'$. By the choice of $H'$, for each triple  in $T_v$, one of these common neighbours 
sees none of the triple.  Hence again,  there is a choice of edges between $P$ and the triple  which yields that $F_i \cap F_{i^*}-H-H'$
contains  a $P_6$. Thus, the number of choices of edges between the union of the elements of $T_v$ and the union of the elements 
of the $P_v$ is at most $2^{6|T_v||P_v|}(1-\frac{1}{2^6})^{|T_v||P_v|}= O( 2^{6|T_v||P_v|- n^{\frac{1}{7(\alpha(T)-1)}}|K|})$. 
If $|K|<20$ then again we have that there are at least $\frac{n}{10(\alpha(T)-1)}$ vertices in $F_i-H-H'-K$ and this set which forms the neighbourhood of $v$ 
in $F_i-H-H'$ is also the common neighbourhood of the vertices of $K$, and hence by the choice of $H'$, for any triple in $T_v$, one of these common neighbours 
is nonadjacent to all of the triple. 
Thus, again the choices for the edges between $K$ and $F_{i^*}-H-H'$ is $2^{|K|(|\pi_{i^*}|-H-H'|-\Omega(n^\frac{1}{7(\alpha(T)-1)})}$. 
On the other hand,  since $K$ is the disjoint union of a clique and a stable set, the number of choices for $K$ and the edges from its vertices within $F_i-H-H'$  is at most ${n \choose k}n 2^{|K|}=O(2^{(2|K|+1)\log n})$. 

It  follows that given  the choice for the pattern 
on $F_{i^*}-H-H'$, there are $2^{O(n)}$ choices for the rest of the pattern. By the maximality of $H$,  the components  of $\overline{F_{i^*}-H}$
are the disjoint union of a clique and a stable set. Mimicking an earlier argument we obtain that the proportion of $T$-free graphs extending patterns we are considering where 
the number $k$ of  vertices in components of $F_{i^*}-H$ which are singletons or have size exceeding $(\log n)^2$  exceeds  $\frac{n}{\sqrt{\log n}}$ is a $o(1)$ proportion 
of the $T$-free graphs.  Hence we can and do focus on patterns for which this is not the case. 

We now forget about $H'$ and construct an $H''$ by  repeating the above argument replacing $C$ by $F_{i^*}-H$. 


For each component $J$ of $F_{i^*}-H-H''$ we let $v_j$ be a vertex which in $G$ sees none of $J-v_J$. 
For $v\in V-F_{i^*}-H-H''$ we define $J'(v)$ to be the set of vertices of $J$ which are adjacent to $v$ if $v_J$ is not, or are nonadjacent to $v$ otherwise. 

We define a set $H''$ by repeatedly adding to it  if possible either 
\begin{enumerate}
\item[(a)]  a vertex which  is  in $F_i-H-H''$ for some $i \neq {i^*}$  and such that either $v$ is nonadajacent to fewer than $\frac{n}{10(\alpha(T)-1)}$
vertices of $\pi_{i^*}-H-H''$ or the sum over all the components  of  $\overline{F_{i^*}-H-H''} $ of  $|J'(v)|$
is less than $\frac{n}{10(\alpha(T)-1)}$, or 
\item[(b)] a triple of vertices   in $F_{i^*}-H-H''$  such that  for some $i \neq {i^*}$  there is a vertex $u$ of $F_i-H-H''$ such that $u$  has at least
$\frac{n}{10(\alpha(T)-1)}$ neighbours in $F_i-H-H''$ and each such neighbour has an edge to the triple.  
\end{enumerate}

We note that for some $\psi>0$. depending only on $T$, the number of choices for the vertices of $H''$ and the edges from them is $O(2^{bn}2^{|H''|(1-\frac{1}{\alpha(T)-1}-\psi)n})$.
So, for all $T$ there is a $C'_T$ such that  the subset of the   $T$-free graphs  we are considering for which $|H''|>C'_T \log n$ is a $o(1)$ proportion of the $T$-free graphs. 
Hence we can  and do restrict our attention to $G$ for which we stop growing $H''$ before it contains $C'_T \log n$ vertices. 


We note that at this point, for every $v \in V-F_{i^*}-H-H''$  there  is a set $T_v$ of at least $\frac{n}{(\log n)^3}$ disjoint triples  in $F_{i^*}-H-H''$ which together with $v$ induce a $P_4$, in which $v$ is an endpoint. 

By the maximality of $H$,  for every $v \in F_i-H-H''$, none of the $P_4$s  induced by $v$ and a triple in $T_v$ extend to an  $P_6$ in $F_{i^*} \cup F_i-H$.
For each $i \neq i^*$,  for each component $K$ of $F_i-H-H''$ which has size at least 2. 
We apply  this fact to a vertex  $v$ in $K$ missing the rest of $K$. 


If $|K| >20$, we greedily  select a family $P_v$   of $\lceil \frac{|K|}{20} \rceil$ pairs 
preferring pairs of adjacent vertices. For each  adjacent pair in $P_v$ and triple in $T_v$ there is a choice of edges between them which yields that together 
with $v$ they yield a $P_6$. If $P$ is a pair of nonadjacent vertices in $P_v$ then there is a stable set in $K$ containing all but  less than $\frac{|K|+1}{10}$
vertices. Hence $|K|<\frac{4n}{5(\alpha(T)-1)}$. Thus there are at least $\frac{n}{10(\alpha(T)-1)}$ vertices in $F_i-H-H''-K$ and this set which forms the neighbourhood of $v$ 
in $F_i-H-H''$ is also the common neighbourhood of $v$ and the vertices in $P$ in $F_i-H-H''$. By the choice of $H''$, for each triple  in $T_v$, one of these common neighbours 
sees none of the triple.  Hence again,  there is a choice of edges between $P$ and the triple  which yields that $F_i \cap F_{i^*}-H-H''$
contains  a $P_6$. Thus, the number of choices of edges between the union of the elements of $T_v$ and the union of the elements 
of the $P_v$ is at most $2^{6|T_v||P_v|}(1-\frac{1}{2^6})^{|T_v||P_v|}=  2^{6|T_v||P_v|- \Omega(\frac{|K|n)}{(\log n)^3})}$. 
If $|K|<20$ then again we have that there are at least $\frac{n}{10(\alpha(T)-1)}$ vertices in $F_i-H-H''-K$ and this set which forms the neighbourhood of $v$ 
in $F_i-H-H''$ is also the common neighbourhood of the vertices of $K$, and hence by the choice of $H'$, for any triple in $T_v$, one of these common neighbours 
is nonadjacent to all of the triple. 
Thus, again the choices for the edges between $K$ and $F_{i^*}-H-H''$ is $2^{|K|(|\pi_i{|^*}-H-H'|-\Omega(\frac{n}{(\log n)^3}))}$. 
On the other hand,  since $K-v$ is the disjoint union of a clique and a stable set, , the number of choices for $K$ and the edges from its vertices within $F_i-H-H''$  is at most ${n \choose |K|}2^{|K|}
=O(2^{(2|K|+1)\log n})$. 

It  follows that the subset of the $T$-free graphs extending the patterns we are considering for which $|H''|>(\log n)^5$ is a $o(1)$ proportion of the $T$-free graphs.
Hence we can and do restrict ourselves to patterns for which this is not the case. We note that  the restriction of $\pi$ to $G-H-H''$ is a $T$-freeness witnessing partition 
where for $i \neq i^*$, $\pi_i-H-H''$ is a clique. 

We set $t_i(v)=m\in((F_i \cap N(v), F_i-N(v))$ and call $v$ extreme on $\pi_i$ if for some sufficiently small $\psi>0$,  $t_i(v)< \psi n$. 

We consider first the possibility that for some $v$, $\min_{i=1}^{\alpha(T)-1} t_i(v)> \frac{n}{\log \log n}$. 
Our bound on the number of  vertices in singleton or large components   of $\overline {F_{i^*}-H}$, implies   $F_i-H$ contains 
 $\omega(\frac{n}{(\log n)^4})$ disjoint triples with  all of which  $v$ induces   $J$ which is either a $P_4$ or the disjoint union of a $P_3$ and a vertex. 
Further since for $i \neq i^*$, $F_i-H-H''$ is a clique we obtain that  $F_i$,contains 
 a family of  $\frac{n}{3 (\log \log n)}$ disjoint edges for  each of which which $v$ sees one endpoint and 
a family of  $\frac{n}{3 (\log \log n)}$ disjoint  edges for  each of which which $v$ sees no  endpoint.
 It follows that  for such a choice of edges, $v$ is
$n^{2/3}$-choice destroying  for the partition $v,F_1-H,F_2-H,...,F_{\alpha(T)-1}-H$. 
So by Lemma \ref{dangerous2cor} , the  number  of $T$-free graphs  extending the type of pattern we are considering  for which such a $v$ exists are  at most 
$2^n2^{-\omega(n^{4/3})}2^{(1-\frac{1}{\alpha(T)-1})\frac{n^2}{2}}$. Hence they are     an $o(1)$ proportion of the $T$-free graphs on $n$ vertices. Thus, we need only consider choices for which no such $v$ exists and every vertex is extreme on at least one $F_i-H_i$. 

We  let $DoublyExtreme$ be the subset of $H \cup H''$ which are extreme on $\pi_i$ for two $i$. 
For any $v$ added to DoublyExtreme, the number of choices for the edges from 
$v$ in $G$ is at most $2^{(1-\frac{2}{\alpha(T)-1})n+2\psi n}$. 

We partition $H -DoublyExtreme$ into $Extreme_1,...,Extreme_{\alpha(T)-1}$ where the vertices of $Extreme_i$ are extreme on $\pi_i$.
We consider the partition $ \pi'$  of $V-DoublyExtreme$ where $\pi'_i=(F_i-H) \cup Extreme_i$. 
We note that, by our bound on $k$, for $v \in  (F_i -H) \cup Extreme_i$, $t_i(v)<\frac{n}{\log \log n}$.

We recursively construct $H'''$ by adding in  (a)  sets of  four vertices  in  $\pi'_{i^*} -H'''$ which induce either 
a $P_4$ or a graph with 2 edges, or (b) sets of six vertices in $F_{i^*} \cup F_i-H$  for some $i \neq i^*$ which induce a $P_6$.
We note that the restriction of $\pi'_i$ to  $G-(H \cup H'' -Extreme_{i^*})-H'''-DoublyExtreme$ is a $T$-freeness certifying partition. 
We note that  the proportion of $T$-free graphs which are extensions of patterns we are considering for which  for any set $S$ of size 4 added to $H'''$, the choices of edges from $S$ make $S$ $n^{2/3}$-choice destroying for this partition  or for any set of size 6  added to $H'''$, the choices of edges from $S$ make $S$ $n^{2/3}$-choice destroying for the coarsening of the partition where we combine the two elements intersecting $S$  is  a $o(1)$ proportion of the $T$-free graphs. Hence we can and do assume there is no such $S$. 
Hence for each such $S$,  for some $\mu>0$, the number of choices for $S$ and  the edges from it is at most $n^{|S|}{n \choose n/ \log \log n}^{|S|}(\alpha(T)-1)2^{n(1-\frac{1}{\alpha(T)-1}-\mu)}$.
Hence the number of choices for $H'''$ and the edges from it is at most $2^{|H'''|n(1-\frac{1}{\alpha(T)-1}-\frac{\mu}{8})}$


We let $t^*=\max_{j\neq {i^*}} (\max _{v \in Extreme_j } t_j(v))$ and let $v^*$ and $j^*$ be such that $v^* \in Extreme_{j^*}$ and $t_{j^*}(v^*)=t^*$.
There are   $l \ge \frac{n}{(\log n)^5}$ triples in $F_{i^*}-H'''-DoublyExtreme$ with which $v$ induces a $P_4$ or the disjoint union of a $P_3$ and a vertex.

If $t^* > (\log n)^9$,  then there are at least $\frac{t^*}{4}$ edges of $F'_{j^*}-H'''-DoublyExtreme$ such that for each choice of a triple  and an edge $e$,
some choice of the edges between the triple and the edge yields a $P_6$ induced by the triple, the edge and $v^*$. By the maximality of $H'''$,
the proportion of choices for the edges between   $F'_{j^*}-H'''-DoublyExtreme$ and $F'_{i^*}-H'''-DoublyExtreme$ which extend such a choice of pattern   is at most  $(1-\frac{1}{2^6})^{t^*l}<e^{-\frac{nt^*}{(\log n)^6}}$.
On the other hand the number of choice of the edges  within the patterns from the vertices in the $Extreme_i$ for $i \neq j$ 
is ${n \choose |H  \cup H''|}{n \choose t^*}^{|H \cup H''|}=e^{O(t^*(\log n)^7}$. So, in this case we are done. 

If $t^*<(\log n)^9$ the number of choice of the edges  within the patterns from the vertices in the $Extreme_i$ for $i \neq j$ 
is $e^{o((\log n)^{16})}$. In this case, if  some $F'_i-H''-DoublyExtreme$ is not a clique then by the maximality of $H$ the  proportion of choices for the edges between   $F'_i-H'''-DoublyExtreme$ and $F'_{i^*}-H'''-DoublyExtreme$ which extend such a choice of pattern to a $T$-free graph  is at most  $(1-\frac{1}{2^6})^{l}<e^{-\frac{n}{(\log n)^6)}}$ and again we are done. So,
$\pi'$ is a $T$-freeness certifying partition of $G-H'''-DoublyExtreme$. Hence if $DoublyExtreme \cup H'''$ is nonempty, $G$ has a $T$-freeness certifying partition,
and there are no bad graphs extending it. Otherwise, we are done by our bounds on the number of edges leaving $H'''$ and $DoublyExtreme$. 
\end{proof}

\end{document}
